\section{Introduction}
Autonomous rendezvous and docking (ARD) is a key enabling technology for a wide range of missions including on-orbit servicing, inspection, refueling, assembly, and active debris removal \cite{fehse2003rendezvous,flores2014}. During a close-proximity operation, the chaser must determine the relative state of the target and use that information to safely guide its own translational and rotational motion towards the docking configuration. Thus, reliable relative navigation is a fundamental prerequisite for autonomous rendezvous and proximity operations (RPO). Vision-based navigation is attractive for such operations because cameras provide rich geometric information at low mass, volume, and power compared with active ranging sensors. Previous studies have demonstrated vision-based relative pose estimation, visual navigation testbeds, and autonomous rendezvous with non-cooperative targets using onboard vision systems \cite{kelsey2006,testbed2014,gaias2018avanti}.

A challenging aspect for vision-based rendezvous is that errors in the perception subsystem can lead to exponentially large errors in the guidance-and-control (GNC) stages. The image plane measurement noises, target motion, incorrect feature correspondences, and imperfect geometric meshes can introduce significant errors to the system. These errors then affect the controller and degrade the safety and success of the rendezvous maneuver. 

For close range pose estimation, the Perspective-$n$-point (PnP) problem provides a direct means of recovering 6 degrees of freedom relative pose from the known three dimensional target points and their two dimensional image projections. The Efficient PnP method provides an efficient $0(n)$ solution by representing the target points using four virtual control points \cite{lepetit2009epnp}. However PnP can be sensitive to noises and thus to develop that we are using Random Sample Consensus (RANSAC) which provides a general robust estimation \cite{fischler1981ransac}. Nonlinear refinement can then be used to minimise reprojection error and improve initial pose estimation.

The relative translational dynamics of a chaser spacecraft near a target in a near-circular LEO are well described by the Hill--Clohessy--Wiltshire (HCW) linearized equations \cite{clohessy1960}. 

This paper presents a modular simulation framework for vision-based
autonomous spacecraft rendezvous with a moving and potentially tumbling target.
The framework integrates relative orbital dynamics, target attitude dynamics,
synthetic camera measurement, pose estimation, nonlinear state estimation, and
constrained guidance and control within a single closed-loop environment, so
that errors and measurement uncertainties propagate through the navigation and
control pipeline exactly as they would onboard.

What distinguishes this framework from the testbeds and benchmarks reviewed in
Section~\ref{sec:related} is that the loop is causally closed. Pose-estimation
benchmarks score a front end on a fixed image set, so the camera trajectory
that produced those images was not generated by a controller acting on the
front end's own output; guidance studies, conversely, assume the state
estimate arrives at the control rate. Here the camera viewpoint at step $k+1$
is a consequence of the pose solution at step $k$, and a lost measurement
therefore propagates. That coupling is what makes the paper's central quantity
observable at all: not how accurate the pose solution is, but how often it
exists. Measurement availability is not a free parameter of the perception
system --- it is a geometric consequence of the same approach trajectory the
controller is generating, and it degrades precisely in the terminal phase,
where the control authority available to recover is smallest.

The primary contributions of this work are as follows:

\begin{itemize}
    \item A closed-loop simulation architecture in which relative dynamics,
    target attitude motion, synthetic vision, pose estimation, state
    estimation, and guidance and control are causally coupled, so that
    measurement dropout is generated by the trajectory rather than imposed
    on it.

    \item The identification of measurement \emph{availability}, rather than
    pose accuracy, as the binding constraint on the closed loop, together
    with a closed-form criterion --- the field-of-view fill range
    $z_{fill} = fL/W$ --- that predicts the range at which availability
    collapses, validated against the measured knee.

    \item A dispersed Monte Carlo characterisation, over randomised initial
    range, bearing, closing speed, tumble rate, target attitude and
    navigation initialisation error, showing that closed-loop performance is
    insensitive to pose accuracy across an order of magnitude of image-plane
    noise and then fails by a threshold mechanism, with the failure
    attributable to availability rather than accuracy.

    \item A systematic evaluation of EPnP-based relative pose estimation under
    increasing image-plane noise, including the planar control-point
    degeneracy and the effect of refining the control-point coefficients
    before nonlinear refinement of the pose itself.

    \item An evaluation of RANSAC-based correspondence rejection under
    controlled outlier contamination, reporting both pose accuracy and
    pose-estimation failure rate.

    \item A numerical validation of the analytical HCW state-transition
    implementation against high-accuracy DOP853 integration, with a
    convergence study in solver tolerance, establishing the dynamics as a
    trustworthy reference for the subsequent closed-loop results.
\end{itemize}
